% ---
% id: meat-core/rem:plus2
% kind: rem
% title: "The constant in the key length and the source's convention"
% summary: >-
%   The key-length bound carries the constant +1 where its source derivation
%   and the published fixed-length statement carry +2: one announced bit buys
%   secrecy on the accepted state by an elementary contraction. The source's
%   secrecy definition, conditioned on the acceptance test alone, does not
%   spend it.
% prominence: standard
% status: accepted
% uses: [thm:keylength]
% mentions: [req:secrecy-convention]
% assumes: []
% refs:
%   - "kamin2025renyi, Thm. 7, Eq. (44)"
% provenance: newly-derived
% lean:
% ---

\begin{remark}
\label{rem:plus2}
The derivation this module's key-length theorem generalises proves
\eqref{eq:secrecy-at} with $\lambda=\lEC+\lceil\log_2\frac1{\eev}\rceil$ and
stops there, i.e.\ its secrecy statement is for the $\wedge\OmAT$ state and
its constant is $+2$; the published fixed-length statement of
\cite[Thm.~7, Eq.~(44)]{kamin2025renyi} reaches the accepted state by a
different route (conditioning on the verification event), with the same
constant. This module spends one deterministic bit (the announced flag) to
reach the accepted state by an elementary contraction argument instead,
avoiding any further import. The cost is exactly one bit in
\eqref{eq:keylength}, which the $\wedge\OmAT$ secrecy definition of that
derivation does not spend; what an instance adopting that definition declares
is Requirement~\ref{req:secrecy-convention}.
\end{remark}
